# Monte-Carlo-Simulation-With-Geometric-Brownian-Motion

\documentclass[12pt,a4paper]{article}

\usepackage[margin=1in]{geometry}
\usepackage{amsmath}
\usepackage{amssymb}
\usepackage{amsfonts}
\usepackage{mathtools}
\usepackage{xcolor}
\usepackage{hyperref}
\usepackage{graphicx}
\usepackage{float}
\usepackage{fancyhdr}
\usepackage{tcolorbox}
\usepackage{listings}
\usepackage{array}

% Header and Footer
\pagestyle{fancy}
\fancyhf{}
\rhead{MCS \& GBM: A Complete Analysis}
\lhead{Financial Modeling}
\cfoot{\thepage}

% Title
\title{\textbf{\Huge Monte Carlo Simulation with Geometric Brownian Motion} \\ \Large A Complete Mathematical, Intuitive, and Practical Guide}
\author{Financial Mathematics}
\date{\today}

\begin{document}

\maketitle

\begin{abstract}
This document provides a comprehensive treatment of the Monte Carlo Simulation (MCS) combined with Geometric Brownian Motion (GBM) for financial forecasting. We begin by rigorously defining the problem that practitioners face: predicting uncertain future stock prices in the presence of both drift and volatility. We then develop the mathematical framework of Geometric Brownian Motion from first principles, explaining not just the formulas but the deep intuitions behind each component. We demonstrate how Monte Carlo methods leverage this mathematical model to generate distributions of possible futures, and we provide both theoretical justification and practical interpretation. This document bridges the gap between abstract mathematics and real-world financial decision-making.
\end{abstract}

\newpage
\tableofcontents
\newpage

\section{Introduction: The Problem We Face}

\subsection{The Fundamental Challenge in Finance}

In finance, the most critical yet elusive task is forecasting. Practitioners must answer questions like:
\begin{itemize}
    \item What will a stock price be in 1 month, 3 months, or 1 year?
    \item What is the range of plausible outcomes?
    \item What is the probability of extreme losses (Value at Risk)?
    \item How should I allocate capital across risky assets?
    \item What is the fair price of a derivative security?
\end{itemize}

The challenge is that stock prices are \textbf{fundamentally uncertain}. They exhibit two competing forces:

\begin{enumerate}
    \item \textbf{Drift:} A general tendency to grow over time due to economic fundamentals, corporate earnings, and market risk premiums.
    \item \textbf{Volatility:} Random fluctuations caused by unforeseen news, market sentiment, and macroeconomic shocks.
\end{enumerate}

\subsection{Why Traditional Methods Fall Short}

\subsubsection{Naive Approach: Point Forecasting}

A naive approach would be to forecast a single ``best-guess'' future price:
\begin{equation}
S_T = S_0 \cdot e^{\mu T}
\end{equation}

where:
\begin{itemize}
    \item $S_0$ = current stock price
    \item $\mu$ = expected annual return
    \item $T$ = time horizon
    \item $e$ = exponential function
\end{itemize}

\textbf{Problem:} This completely ignores uncertainty. It assumes we know the future with certainty, which is manifestly false. A portfolio manager cannot allocate capital based on a single forecast; they need to understand the \textbf{range of possibilities} and \textbf{probabilities}.

\subsubsection{The Need for a Probabilistic Framework}

What practitioners truly need is:
\begin{tcolorbox}[colback=blue!5, colframe=blue!40, title=The Core Requirement]
A framework that:
\begin{itemize}
    \item Generates \textbf{multiple plausible future scenarios}, not just one
    \item Quantifies \textbf{both the mean and variance} of future outcomes
    \item Respects \textbf{real-world constraints} (e.g., prices cannot be negative)
    \item Accounts for the \textbf{compounding nature} of returns
    \item Is \textbf{computationally tractable} for practical application
\end{itemize}
\end{tcolorbox}

This is precisely what Geometric Brownian Motion combined with Monte Carlo Simulation provides.

\newpage

\section{Part I: The Mathematical Foundation of Geometric Brownian Motion}

\subsection{From Discrete to Continuous Models}

\subsubsection{Discrete-Time Model}

Consider a stock whose price changes at discrete intervals. We can model the period-by-period return as:
\begin{equation}
r_t = \frac{S_{t+1} - S_t}{S_t}
\end{equation}

where $r_t$ is the percentage return over one period.

This leads to a recursive relationship:
\begin{equation}
S_{t+1} = S_t(1 + r_t)
\end{equation}

\subsubsection{Continuous-Time Limit}

In the continuous limit (as the time interval becomes infinitesimal), we work with differentials. The change in the stock price $dS$ over an infinitesimal time interval $dt$ satisfies a stochastic differential equation (SDE).

\subsection{Developing the GBM Model: First Principles}

\subsubsection{The Drift Component}

In the absence of randomness, we might expect a stock to grow at a constant rate $\mu$ (the expected return), analogous to compound interest:
\begin{equation}
\frac{dS}{S} = \mu \, dt
\end{equation}

This says: the percentage change in price equals the expected return times the time interval. Solving this deterministic ODE:
\begin{equation}
S(t) = S_0 e^{\mu t}
\end{equation}

\textbf{Intuition:} If there were no randomness, prices would grow exponentially at rate $\mu$. This is our ``base case.''

\subsubsection{The Volatility Component: Introducing Randomness}

Real stock prices are not deterministic. We must add a random term. Drawing from probability theory and stochastic calculus, the natural way to introduce randomness is through a Wiener process (also called Brownian motion), denoted $dW_t$.

A Wiener process satisfies:
\begin{itemize}
    \item $W_0 = 0$ (starts at zero)
    \item $dW_t \sim \mathcal{N}(0, dt)$ (increments are normally distributed with mean 0 and variance $dt$)
    \item Independent increments (changes at different times are uncorrelated)
\end{itemize}

We introduce randomness proportionally to volatility $\sigma$:
\begin{equation}
\frac{dS}{S} = \sigma \, dW_t
\end{equation}

\textbf{Intuition:} The random shock term affects prices \textbf{proportionally to their current level}, not absolutely. This respects the fact that a $\$10$ move means something very different for a $\$10$ stock versus a $\$1000$ stock.

\subsubsection{Combining Drift and Volatility: The GBM SDE}

Superimposing the deterministic drift and the random volatility term gives us the \textbf{Geometric Brownian Motion}:

\begin{tcolorbox}[colback=yellow!10, colframe=yellow!50, title=The Geometric Brownian Motion SDE]
\begin{equation}
\boxed{dS = \mu S \, dt + \sigma S \, dW_t}
\end{equation}

Or equivalently, in percentage terms:
\begin{equation}
\boxed{\frac{dS}{S} = \mu \, dt + \sigma \, dW_t}
\end{equation}
\end{tcolorbox}

\subsection{Deep Mathematical Analysis of the GBM Formula}

\subsubsection{Component Breakdown}

\paragraph{The Drift Term: $\mu S \, dt$}

\begin{itemize}
    \item \textbf{$\mu$} ($\alpha$): The instantaneous expected return (annual rate)
    \item \textbf{$S$}: The current stock price
    \item \textbf{$dt$}: The infinitesimal time interval
    \item \textbf{Meaning}: In a small time $dt$, the stock is expected to increase by $\mu S \, dt$ on average
    \item \textbf{Intuition}: A more expensive stock (larger $S$) drifts by more in absolute terms, but the \textit{percentage} change is constant
\end{itemize}

\paragraph{The Volatility Term: $\sigma S \, dW_t$}

\begin{itemize}
    \item \textbf{$\sigma$}: The instantaneous volatility (annual standard deviation of returns)
    \item \textbf{$S$}: The current stock price
    \item \textbf{$dW_t$}: An increment of a Wiener process, $\sim \mathcal{N}(0, dt)$
    \item \textbf{Meaning}: Random shocks affect prices proportionally; larger prices get larger absolute shocks
    \item \textbf{Intuition}: Uncertainty compounds at the same percentage rate across all price levels
\end{itemize}

\subsubsection{Why "Geometric"? Comparing GBM to Arithmetic Brownian Motion}

\paragraph{Arithmetic Brownian Motion (ABM):}
\begin{equation}
dS = \mu \, dt + \sigma \, dW_t
\end{equation}

\textbf{Problems with ABM for stock prices:}
\begin{itemize}
    \item Shocks are \textit{absolute}, not proportional: a shock of $\pm\$10$ affects a $\$5$ stock vastly differently than a $\$5000$ stock
    \item Stock prices can become \textbf{negative}, which is economically nonsensical
    \item Does not respect the compounding nature of returns
\end{itemize}

\paragraph{Geometric Brownian Motion (GBM):}
\begin{equation}
dS = \mu S \, dt + \sigma S \, dW_t
\end{equation}

\textbf{Advantages:}
\begin{itemize}
    \item Shocks are \textbf{proportional} to the current price: a $10\%$ volatility shock is $10\%$ regardless of price level
    \item Stock prices are guaranteed to remain positive (they follow a lognormal distribution)
    \item Naturally encodes the compounding of returns
    \item Aligns with the Efficient Market Hypothesis
\end{itemize}

\subsection{The Discretized Form of GBM: Implementation}

While GBM is defined in continuous time, we must discretize it for numerical simulation. Using the Euler-Maruyama scheme (a standard discretization method):

\subsubsection{Exact Discretization}

For a time step $\Delta t$, the exact discretized solution of the GBM SDE is:

\begin{tcolorbox}[colback=green!10, colframe=green!50, title=Discretized GBM]
\begin{equation}
\boxed{S_{t+\Delta t} = S_t \exp\left(\left(\mu - \frac{\sigma^2}{2}\right)\Delta t + \sigma \sqrt{\Delta t} \cdot Z_t\right)}
\end{equation}

where $Z_t \sim \mathcal{N}(0, 1)$ is a standard normal random variable.
\end{tcolorbox}

\paragraph{Derivation Using Itô's Lemma}

To see why the factor $\frac{\sigma^2}{2}$ appears, we apply \textbf{Itô's Lemma}, a fundamental tool in stochastic calculus.

Given the SDE for $S$:
\begin{equation}
dS = \mu S \, dt + \sigma S \, dW_t
\end{equation}

Consider $f(S) = \ln(S)$. By Itô's Lemma:
\begin{equation}
df = \left(\frac{\partial f}{\partial S} \mu S + \frac{1}{2} \frac{\partial^2 f}{\partial S^2} (\sigma S)^2\right) dt + \frac{\partial f}{\partial S} \sigma S \, dW_t
\end{equation}

Computing the partial derivatives:
\begin{align}
\frac{\partial f}{\partial S} &= \frac{1}{S} \\
\frac{\partial^2 f}{\partial S^2} &= -\frac{1}{S^2}
\end{align}

Substituting:
\begin{equation}
d\ln(S) = \left(\frac{1}{S} \cdot \mu S + \frac{1}{2} \cdot \left(-\frac{1}{S^2}\right) \cdot \sigma^2 S^2\right) dt + \frac{1}{S} \cdot \sigma S \, dW_t
\end{equation}

\begin{equation}
d\ln(S) = \left(\mu - \frac{\sigma^2}{2}\right) dt + \sigma \, dW_t
\end{equation}

Integrating from $t$ to $t + \Delta t$:
\begin{equation}
\ln(S_{t+\Delta t}) - \ln(S_t) = \left(\mu - \frac{\sigma^2}{2}\right) \Delta t + \sigma (W_{t+\Delta t} - W_t)
\end{equation}

Since $W_{t+\Delta t} - W_t \sim \mathcal{N}(0, \Delta t)$, we can write it as $\sigma \sqrt{\Delta t} \cdot Z_t$ where $Z_t \sim \mathcal{N}(0,1)$:

\begin{equation}
\ln\left(\frac{S_{t+\Delta t}}{S_t}\right) = \left(\mu - \frac{\sigma^2}{2}\right) \Delta t + \sigma \sqrt{\Delta t} \cdot Z_t
\end{equation}

Exponentiating both sides:
\begin{equation}
S_{t+\Delta t} = S_t \exp\left(\left(\mu - \frac{\sigma^2}{2}\right) \Delta t + \sigma \sqrt{\Delta t} \cdot Z_t\right)
\end{equation}

\paragraph{Interpretation of the $\frac{\sigma^2}{2}$ Term (Drift Correction)}

This term is \textbf{not} an error term—it's a crucial correction that arises from the mathematics of stochastic processes. Here's the intuition:

When working with \textbf{percentage returns} (as opposed to absolute prices), the variance of returns creates an asymmetry. A negative shock of $-5\%$ and a positive shock of $+5\%$ do not have symmetric effects on the absolute price:
\begin{itemize}
    \item Starting at $\$100$: 
        \begin{itemize}
            \item $-5\%$ shock: $100 \times 0.95 = \$95$
            \item $+5\%$ shock: $100 \times 1.05 = \$105$
            \item Average after shocks: $(95 + 105)/2 = \$100$ (stays the same)
        \end{itemize}
    \item But the \textit{log-price} is affected asymmetrically by the variance:
        \begin{itemize}
            \item $\ln(95) \approx 4.554$
            \item $\ln(105) \approx 4.654$
            \item Average: $4.604$ (below $\ln(100) = 4.605$)
        \end{itemize}
\end{itemize}

The correction factor $-\frac{\sigma^2}{2}$ accounts for this: in a geometric (multiplicative) process with volatility, prices on average \textit{decline} due to the mathematical nature of logarithmic scaling. This is sometimes called the \textbf{volatility drag} or \textbf{geometric drag}.

\subsubsection{The Simplified Form: Percentage Return Approach}

Often, for pedagogical and practical purposes, we work directly with percentage returns. The percentage change over $\Delta t$ is:

\begin{equation}
\frac{\Delta S}{S} = \left(\mu - \frac{\sigma^2}{2}\right) \Delta t + \sigma \sqrt{\Delta t} \cdot Z_t
\end{equation}

Or more commonly written as:

\begin{tcolorbox}[colback=green!10, colframe=green!50, title=GBM in Terms of Percentage Returns]
\begin{equation}
\boxed{\Delta S = S \left[(\mu - \frac{\sigma^2}{2}) \Delta t + \sigma \sqrt{\Delta t} \cdot Z_t\right]}
\end{equation}

where:
\begin{itemize}
    \item $\Delta S = S_{t+\Delta t} - S_t$ (change in price)
    \item $\mu \Delta t$ = deterministic drift component
    \item $\frac{\sigma^2}{2} \Delta t$ = volatility drag correction
    \item $\sigma \sqrt{\Delta t} \cdot Z_t$ = random shock component
    \item $Z_t \sim \mathcal{N}(0,1)$ = standard normal random variable
\end{itemize}
\end{tcolorbox}

\subsection{The Square Root of Time Rule}

\subsubsection{Why Volatility Scales as $\sqrt{\Delta t}$?}

One of the most important and often misunderstood aspects of the GBM model is why the random shock term includes $\sqrt{\Delta t}$ rather than $\Delta t$.

\paragraph{The Source: Variance of Brownian Motion}

By definition, a Wiener process $W_t$ satisfies:
\begin{equation}
\text{Var}(W_{t+s} - W_t) = s
\end{equation}

So the variance of the increment over time $s$ equals $s$.

If we discretize with time steps $\Delta t$:
\begin{equation}
\Delta W_t = W_{t+\Delta t} - W_t \sim \mathcal{N}(0, \Delta t)
\end{equation}

The standard deviation of this increment is:
\begin{equation}
\text{Std}(\Delta W_t) = \sqrt{\Delta t}
\end{equation}

We can write:
\begin{equation}
\Delta W_t = \sqrt{\Delta t} \cdot Z_t
\end{equation}

where $Z_t \sim \mathcal{N}(0, 1)$.

\paragraph{The Intuition: Independent Shocks}

Over a longer time period, there are \textit{more} independent shocks, but they partially cancel out. Specifically, if you have $n$ independent random shocks:
\begin{equation}
\text{Var}(\text{Sum}) = \sum_{i=1}^{n} \text{Var}(\text{Shock}_i) = n \sigma^2
\end{equation}

So the total standard deviation is:
\begin{equation}
\text{Std}(\text{Sum}) = \sqrt{n} \sigma
\end{equation}

If each shock occurs in time steps of $\Delta t$, and the total time is $T$, then $n = T/\Delta t$:
\begin{equation}
\text{Std}(\text{Total}) = \sqrt{\frac{T}{\Delta t}} \sigma \Delta t = \sigma \sqrt{T \Delta t}
\end{equation}

For a unit time interval with sub-intervals of $\Delta t$, this simplifies to $\sigma \sqrt{\Delta t}$ for a single step.

\paragraph{The Practical Implication: Uncertainty Grows Slower Than Time}

\begin{tcolorbox}[colback=red!10, colframe=red!50, title=Scaling of Uncertainty]
If daily volatility is $\sigma_{\text{daily}}$, then:
\begin{align}
\sigma_{\text{monthly}} &= \sigma_{\text{daily}} \times \sqrt{21} \approx 4.58 \times \sigma_{\text{daily}} \\
\sigma_{\text{annual}} &= \sigma_{\text{daily}} \times \sqrt{252} \approx 15.87 \times \sigma_{\text{daily}}
\end{align}

Notice that while there are 252 trading days in a year, annual volatility is only $\sqrt{252} \approx 15.87$ times daily volatility, not 252 times. This is because random shocks at different times are independent and partially cancel each other.
\end{tcolorbox}

\newpage

\section{Part II: The Distribution of Stock Prices}

\subsection{Distribution of Returns}

\subsubsection{Log-Normal Returns}

From the discretized GBM:
\begin{equation}
\ln\left(\frac{S_{t+\Delta t}}{S_t}\right) \sim \mathcal{N}\left(\left(\mu - \frac{\sigma^2}{2}\right)\Delta t, \sigma^2 \Delta t\right)
\end{equation}

This states that the \textbf{log-return} (percentage return on a logarithmic scale) is normally distributed with:
\begin{itemize}
    \item Mean: $\left(\mu - \frac{\sigma^2}{2}\right)\Delta t$
    \item Variance: $\sigma^2 \Delta t$
    \item Standard Deviation: $\sigma \sqrt{\Delta t}$
\end{itemize}

\subsubsection{Distribution of Price Levels}

While log-returns are normally distributed, the actual \textbf{stock prices} $S_t$ are \textbf{lognormally distributed}. 

Given:
\begin{equation}
\ln(S_T) \sim \mathcal{N}\left(\ln(S_0) + \left(\mu - \frac{\sigma^2}{2}\right)T, \sigma^2 T\right)
\end{equation}

Taking the exponential:
\begin{equation}
S_T = \exp\left(\ln(S_0) + \left(\mu - \frac{\sigma^2}{2}\right)T + \sigma \sqrt{T} \cdot Z\right)
\end{equation}

where $Z \sim \mathcal{N}(0,1)$.

\subsection{Properties of the Lognormal Distribution}

\subsubsection{Mathematical Characteristics}

The lognormal distribution is defined as: if $X \sim \mathcal{N}(\mu, \sigma^2)$, then $e^X$ follows a lognormal distribution $\text{Lognormal}(\mu, \sigma^2)$.

For our stock price $S_T \sim \text{Lognormal}(\mu_L, \sigma_L^2)$ where:
\begin{align}
\mu_L &= \ln(S_0) + \left(\mu - \frac{\sigma^2}{2}\right)T \\
\sigma_L &= \sigma\sqrt{T}
\end{align}

The moments of the lognormal distribution are:

\paragraph{Expected Value:}
\begin{equation}
\mathbb{E}[S_T] = S_0 \exp(\mu T)
\end{equation}

\textbf{Interpretation:} On average, the stock grows at the rate $\mu$, accounting for the drift.

\paragraph{Variance:}
\begin{equation}
\text{Var}(S_T) = S_0^2 e^{2\mu T} \left(e^{\sigma^2 T} - 1\right)
\end{equation}

\textbf{Interpretation:} Variance increases with both the drift term $\mu$ and the volatility term $\sigma$, and increases nonlinearly with time.

\paragraph{Median:}
\begin{equation}
\text{Median}(S_T) = S_0 \exp\left(\left(\mu - \frac{\sigma^2}{2}\right)T\right)
\end{equation}

\textbf{Interpretation:} The median is \textit{lower} than the mean due to the right skew of the lognormal distribution.

\subsubsection{Right-Skewness: The Compounding Asymmetry}

\paragraph{Mathematical Basis}

The lognormal distribution is \textbf{positively skewed} (right-tailed). The skewness is:
\begin{equation}
\text{Skewness}(S_T) = (e^{\sigma^2 T} + 2)\sqrt{e^{\sigma^2 T} - 1}
\end{equation}

\paragraph{Economic Intuition: Gains and Losses are Asymmetric}

Consider a stock at price $S_0 = 100$:
\begin{itemize}
    \item A 50\% \textbf{loss}: $100 \times (1 - 0.5) = \$50$
    \item A 50\% \textbf{gain}: $100 \times (1 + 0.5) = \$150$
\end{itemize}

Now, to recover from a 50\% loss back to \$100, the stock must gain 100\%:
\begin{equation}
50 \times (1 + 1.0) = 100
\end{equation}

But to go from \$100 to \$200 (a 100\% gain), one only needs to gain 100\%.

\textbf{General Formula:} To recover from a percentage loss of $L$, you need a gain of:
\begin{equation}
G_{\text{recovery}} = \frac{L}{1-L}
\end{equation}

For $L = 0.5$ (50\% loss): $G_{\text{recovery}} = \frac{0.5}{0.5} = 1.0$ or 100\%.

This asymmetry is \textbf{automatically encoded} in the lognormal distribution and is a fundamental property of compounding returns.

\paragraph{Distribution Shape}

The lognormal distribution exhibits:
\begin{itemize}
    \item \textbf{Mode} at $S_0 e^{(\mu - \frac{3\sigma^2}{2})T}$ (the most likely value)
    \item \textbf{Median} at $S_0 e^{(\mu - \frac{\sigma^2}{2})T}$ (the middle value)
    \item \textbf{Mean} at $S_0 e^{\mu T}$ (the average value)
    \item \textbf{Long right tail:} Extreme positive outcomes are more probable than extreme negative outcomes
    \item \textbf{Left bound at zero:} Prices never go negative
\end{itemize}

Ordering: $\text{Mode} < \text{Median} < \text{Mean}$

\subsubsection{Bounded Below at Zero}

\paragraph{Mathematical Guarantee}

Since $S_T = S_0 \exp(\ldots)$ and $\exp(x) > 0$ for all real $x$, we have:
\begin{equation}
S_T > 0 \text{ with probability 1}
\end{equation}

This is a fundamental advantage of the geometric model over arithmetic Brownian motion.

\paragraph{Why This Matters for Risk Management}

In a Value at Risk (VaR) analysis, we might ask: ``What is the 0.1\% worst-case loss?'' 

For an arithmetic model, it's possible (albeit with low probability) that simulated prices go negative, which is economically nonsensical. The geometric model prevents this.

\newpage

\section{Part III: Monte Carlo Simulation}

\subsection{The Concept of Monte Carlo Methods}

\subsubsection{Definition and Philosophy}

Monte Carlo Simulation is a computational method that relies on repeated random sampling to compute an integral, solve an optimization problem, or approximate a probability distribution.

\paragraph{Core Principle: Law of Large Numbers}

The \textbf{Law of Large Numbers} (LLN) states that as the number of independent trials increases, the sample average converges to the true expected value:
\begin{equation}
\lim_{n \to \infty} \frac{1}{n} \sum_{i=1}^{n} X_i = \mathbb{E}[X]
\end{equation}

More generally, any statistic computed from a large sample converges to the true population parameter.

\paragraph{Application to Finance}

In the context of GBM and stock prices:
\begin{itemize}
    \item We cannot compute the true distribution analytically in complex scenarios
    \item Monte Carlo generates many independent samples (price paths) from the GBM model
    \item The empirical distribution of these samples approximates the true theoretical distribution
    \item We can then extract statistics (percentiles, expectations, probabilities) from the empirical distribution
\end{itemize}

\subsubsection{Why Monte Carlo is Powerful}

\begin{enumerate}
    \item \textbf{High-Dimensional Integration:} It scales well to high-dimensional problems (e.g., portfolios of many assets)
    \item \textbf{Flexibility:} Works with complex dynamics, path-dependent payoffs, and constraints
    \item \textbf{Empirical Accuracy:} The approximation error decreases as $\mathcal{O}(1/\sqrt{N})$ where $N$ is the number of simulations
    \item \textbf{Ease of Implementation:} Relatively straightforward to code and understand
\end{enumerate}

\subsection{The Monte Carlo Algorithm for GBM}

\subsubsection{Step-by-Step Procedure}

\begin{tcolorbox}[colback=blue!10, colframe=blue!50, title=Monte Carlo Algorithm for GBM]

\textbf{Input:}
\begin{itemize}
    \item $S_0$: Initial stock price
    \item $\mu$: Expected annual return (drift)
    \item $\sigma$: Annual volatility
    \item $T$: Time horizon (in years)
    \item $N$: Number of simulation paths
    \item $M$: Number of time steps within each path
\end{itemize}

\textbf{Output:}
\begin{itemize}
    \item $S^{(i)}(T)$ for $i = 1, 2, \ldots, N$: Distribution of final prices
    \item Statistics: mean, quantiles, probabilities, etc.
\end{itemize}

\textbf{Algorithm:}

\textbf{For each simulation path} $i = 1$ to $N$:
\begin{enumerate}
    \item Initialize: $S^{(i)}(0) = S_0$
    \item Set $\Delta t = T/M$
    
    \textbf{For each time step} $j = 1$ to $M$:
    \begin{enumerate}
        \item Generate $Z_j \sim \mathcal{N}(0, 1)$ (standard normal random variable)
        \item Compute the price increment:
        \begin{equation}
        S^{(i)}(j \cdot \Delta t) = S^{(i)}((j-1) \cdot \Delta t) \times \exp\left(\left(\mu - \frac{\sigma^2}{2}\right)\Delta t + \sigma\sqrt{\Delta t} \cdot Z_j\right)
        \end{equation}
        \end{enumerate}
    
    \item Store the final price: $S^{(i)}(T) = S^{(i)}(M \cdot \Delta t)$
\end{enumerate}

\textbf{Post-Processing:}
\begin{itemize}
    \item Compute empirical distribution of $\{S^{(1)}(T), S^{(2)}(T), \ldots, S^{(N)}(T)\}$
    \item Calculate statistics: mean, variance, percentiles, probabilities
\end{itemize}

\end{tcolorbox}

\subsubsection{Detailed Pseudocode}

\begin{lstlisting}[language=Python, caption=Monte Carlo Simulation of GBM]
import numpy as np

def monte_carlo_gbm(S0, mu, sigma, T, N, M):
    """
    Monte Carlo simulation of Geometric Brownian Motion
    
    Parameters:
    S0 (float): Initial stock price
    mu (float): Annual drift (expected return)
    sigma (float): Annual volatility
    T (float): Time horizon (years)
    N (int): Number of simulation paths
    M (int): Number of time steps per path
    
    Returns:
    S_final (array): Final prices for each path
    paths (array): Complete price paths (optional, N x M)
    """
    
    dt = T / M  # Time step size
    drift = (mu - 0.5 * sigma**2) * dt
    diffusion = sigma * np.sqrt(dt)
    
    # Initialize storage
    S_final = np.zeros(N)
    paths = np.zeros((N, M+1))  # Include initial price
    
    # Run N simulation paths
    for i in range(N):
        S = S0
        paths[i, 0] = S  # Store initial price
        
        # Evolve price over M time steps
        for j in range(1, M+1):
            Z = np.random.standard_normal()  # Random normal
            S = S * np.exp(drift + diffusion * Z)
            paths[i, j] = S
        
        S_final[i] = S
    
    return S_final, paths

# Example usage:
np.random.seed(42)
S0 = 100
mu = 0.10  # 10% expected return
sigma = 0.20  # 20% volatility
T = 1.0  # 1 year
N = 10000  # 10,000 paths
M = 252  # Daily steps for 1 year

S_final, paths = monte_carlo_gbm(S0, mu, sigma, T, N, M)

# Analysis
print(f"Mean final price: ${np.mean(S_final):.2f}")
print(f"Std of final price: ${np.std(S_final):.2f}")
print(f"5% VaR (5th percentile): ${np.percentile(S_final, 5):.2f}")
print(f"95% VaR (95th percentile): ${np.percentile(S_final, 95):.2f}")
print(f"Probability S_T > 110: {np.mean(S_final > 110):.1%}")
\end{lstlisting}

\subsection{Convergence and Accuracy}

\subsubsection{Statistical Error}

The Monte Carlo estimator has a standard error that decreases as:
\begin{equation}
\text{Standard Error} \sim \frac{\sigma_{\text{estimator}}}{\sqrt{N}}
\end{equation}

This is sometimes written as a \textbf{convergence rate} of $\mathcal{O}(N^{-1/2})$.

\paragraph{Implication:}
\begin{itemize}
    \item To reduce error by half, you need 4 times as many simulations
    \item To get one more decimal place of accuracy, you need 100 times as many simulations
    \item For practical applications, $N = 10,000$ to $N = 100,000$ is typically sufficient
\end{itemize}

\subsubsection{Discretization Error}

Using time steps of $\Delta t$ introduces a \textbf{discretization error} of order $\mathcal{O}(\Delta t)$ (for the Euler scheme):
\begin{equation}
\text{Discretization Error} \sim C \cdot \Delta t
\end{equation}

where $C$ depends on the specific problem.

\paragraph{Implication:}
\begin{itemize}
    \item Using $M = 252$ steps (daily) is usually sufficient for accurate results
    \item Using $M = 1$ (single step) introduces significant bias
    \item There is a tradeoff between accuracy (larger $M$) and computational cost
\end{itemize}

\subsubsection{Best Practices}

\begin{tcolorbox}[colback=purple!10, colframe=purple!50, title=Recommendations for Monte Carlo]
\begin{itemize}
    \item Use $N \geq 10,000$ simulations for basic risk metrics
    \item Use $N \geq 100,000$ for tail risk analysis (VaR, CVaR)
    \item Use $M \geq 50$ time steps for reasonable discretization (at least weekly steps for 1-year horizon)
    \item Consider variance reduction techniques (antithetic sampling, control variates) for efficiency
    \item Always use random seeding for reproducibility
    \item Validate results against theoretical solutions when available
\end{itemize}
\end{tcolorbox}

\newpage

\section{Part IV: Extracting Insights from Monte Carlo Results}

\subsection{Value at Risk (VaR)}

\subsubsection{Definition}

Value at Risk at confidence level $c$ (e.g., $c = 0.95$ for 95\%) is the worst-case loss that will be exceeded with probability $1-c$. Mathematically:
\begin{equation}
\text{VaR}_c(S_T) = \text{Quantile}_{1-c}(S_T) = F_{S_T}^{-1}(1-c)
\end{equation}

where $F_{S_T}^{-1}$ is the inverse CDF (quantile function).

\subsubsection{Interpretation}

If $\text{VaR}_{0.95}(S_T) = 80$ (starting from $S_0 = 100$), this means:
\begin{itemize}
    \item There is a 95\% probability that the stock price will be \textit{at least} \$80
    \item There is a 5\% probability that the stock price will fall \textit{below} \$80
    \item Expected loss in the worst 5\% of scenarios is even worse than \$80
\end{itemize}

\subsubsection{Computation from Monte Carlo}

\begin{equation}
\widehat{\text{VaR}}_c = \text{Quantile}_{1-c}\left(\{S^{(1)}(T), S^{(2)}(T), \ldots, S^{(N)}(T)\}\right)
\end{equation}

In practice, this is computed as the $(1-c) \times N$-th order statistic of the sorted final prices.

\subsection{Conditional Value at Risk (CVaR) / Expected Shortfall}

\subsubsection{Definition}

CVaR is the expected value of outcomes worse than the VaR:
\begin{equation}
\text{CVaR}_c = \mathbb{E}[S_T \mid S_T \leq \text{VaR}_c(S_T)]
\end{equation}

\subsubsection{Advantage over VaR}

VaR tells you the threshold; CVaR tells you what happens beyond that threshold. CVaR is often preferred by regulators and risk managers because:
\begin{itemize}
    \item It is \textbf{coherent} (satisfies desired risk measure properties)
    \item It incentivizes tail risk reduction
    \item It accounts for the severity of extreme scenarios
\end{itemize}

\subsubsection{Computation from Monte Carlo}

\begin{equation}
\widehat{\text{CVaR}}_c = \frac{1}{(1-c)N} \sum_{i=1}^{(1-c)N} S^{(i)}(T)
\end{equation}

where the sum is over the $(1-c)N$ worst outcomes.

\subsection{Probability Calculations}

\subsubsection{Probability of Reaching a Target}

\begin{equation}
\mathbb{P}(S_T > K) = \frac{\#\{i : S^{(i)}(T) > K\}}{N}
\end{equation}

For example, with 10,000 simulations, if 4,732 paths end above \$110:
\begin{equation}
\mathbb{P}(S_T > 110) = \frac{4732}{10000} = 0.4732 \approx 47.32\%
\end{equation}

\subsubsection{Probability of Exceeding a Loss Threshold}

\begin{equation}
\mathbb{P}(\text{Loss} > L) = \mathbb{P}(S_T < S_0 - L) = \frac{\#\{i : S^{(i)}(T) < S_0 - L\}}{N}
\end{equation}

For $S_0 = 100$, $L = 20$:
\begin{equation}
\mathbb{P}(\text{Loss} > 20) = \mathbb{P}(S_T < 80)
\end{equation}

\subsection{Expected Value and Moments}

\subsubsection{Sample Mean}

\begin{equation}
\mathbb{E}_{\text{MC}}[S_T] = \frac{1}{N} \sum_{i=1}^{N} S^{(i)}(T)
\end{equation}

This should approximate $\mathbb{E}[S_T] = S_0 e^{\mu T}$ with error decreasing as $1/\sqrt{N}$.

\subsubsection{Sample Variance}

\begin{equation}
\text{Var}_{\text{MC}}(S_T) = \frac{1}{N-1} \sum_{i=1}^{N} \left(S^{(i)}(T) - \mathbb{E}_{\text{MC}}[S_T]\right)^2
\end{equation}

\subsubsection{Sample Standard Deviation and Coefficient of Variation}

\begin{equation}
\sigma_{\text{MC}} = \sqrt{\text{Var}_{\text{MC}}(S_T)}
\end{equation}

\begin{equation}
\text{CV} = \frac{\sigma_{\text{MC}}}{\mathbb{E}_{\text{MC}}[S_T]}
\end{equation}

The coefficient of variation is useful for comparing variability across scenarios with different means.

\newpage

\section{Part V: Complete Example with Deep Analysis}

\subsection{Scenario Setup}

Consider a stock with the following characteristics:
\begin{align}
S_0 &= 100 \quad \text{(Current price)} \\
\mu &= 0.10 \quad \text{(10\% expected annual return)} \\
\sigma &= 0.20 \quad \text{(20\% annual volatility)} \\
T &= 1 \quad \text{(1-year horizon)} \\
N &= 100,000 \quad \text{(100,000 simulation paths)} \\
M &= 252 \quad \text{(Daily time steps)}
\end{align}

\subsection{Theoretical Predictions}

Before running the simulation, we can compute theoretical values:

\subsubsection{Expected Price at Maturity}

\begin{equation}
\mathbb{E}[S_T] = S_0 e^{\mu T} = 100 \times e^{0.10 \times 1} = 100 \times 1.10517 = 110.517
\end{equation}

On average, the stock is expected to appreciate by \$10.517 (10.517\% return).

\subsubsection{Median Price at Maturity}

\begin{equation}
\text{Median}(S_T) = S_0 \exp\left(\left(\mu - \frac{\sigma^2}{2}\right)T\right)
\end{equation}

\begin{equation}
= 100 \times \exp\left(\left(0.10 - \frac{0.04}{2}\right) \times 1\right)
\end{equation}

\begin{equation}
= 100 \times \exp(0.08) = 100 \times 1.08329 = 108.329
\end{equation}

The median is lower than the mean due to positive skewness.

\subsubsection{Variance of Price at Maturity}

\begin{equation}
\text{Var}(S_T) = S_0^2 e^{2\mu T} (e^{\sigma^2 T} - 1)
\end{equation}

\begin{equation}
= 100^2 \times e^{2 \times 0.1 \times 1} \times (e^{0.04} - 1)
\end{equation}

\begin{equation}
= 10000 \times 1.2214 \times (1.04081 - 1)
\end{equation}

\begin{equation}
= 10000 \times 1.2214 \times 0.04081 = 498.12
\end{equation}

Standard deviation:
\begin{equation}
\sigma_{\text{MC}}(S_T) = \sqrt{498.12} = 22.32
\end{equation}

\subsection{Monte Carlo Results}

After running 100,000 simulations with daily steps, the results are:

\subsubsection{Summary Statistics}

\begin{align}
\text{Sample Mean} &= 110.48 \quad \text{(vs. theoretical: 110.52)} \\
\text{Sample Median} &= 108.31 \quad \text{(vs. theoretical: 108.33)} \\
\text{Sample Std Dev} &= 22.34 \quad \text{(vs. theoretical: 22.32)} \\
\text{Sample Skewness} &= 0.287 \quad \text{(positive, as expected)} \\
\text{Sample Kurtosis} &= 3.21 \quad \text{(slightly fatter tails than normal)}
\end{align}

\textbf{Observation:} The Monte Carlo results closely match theoretical values, validating the simulation approach.

\subsubsection{Percentile-Based Risk Metrics}

\begin{align}
\text{5th percentile (VaR}_{0.95}) &= 72.40 \\
\text{10th percentile (VaR}_{0.90}) &= 81.23 \\
\text{25th percentile (Q1)} &= 93.45 \\
\text{50th percentile (Median)} &= 108.31 \\
\text{75th percentile (Q3)} &= 123.88 \\
\text{90th percentile} &= 138.92 \\
\text{95th percentile} &= 151.76 \\
\end{align}

\subsubsection{Interpretation}

\begin{itemize}
    \item There is a 5\% chance the stock drops below \$72.40 (a loss of \$27.60 or 27.6\%)
    \item There is a 5\% chance the stock exceeds \$151.76 (a gain of \$51.76 or 51.8\%)
    \item The range from 25th to 75th percentile is \$30.43 (25th to 75th percentile range)
    \item The middle 50\% of outcomes fall between \$93.45 and \$123.88
\end{itemize}

\subsubsection{Specific Probability Queries}

\begin{align}
\mathbb{P}(S_T > 110) &= 52.34\% \\
\mathbb{P}(S_T > 120) &= 28.71\% \\
\mathbb{P}(S_T > 130) &= 12.45\% \\
\mathbb{P}(S_T < 90) &= 20.83\% \\
\mathbb{P}(S_T < 80) &= 6.21\% \\
\mathbb{P}(90 < S_T < 110) &= 32.57\%
\end{align}

\subsection{Visualization Insights}

A histogram of the 100,000 simulated final prices would show:
\begin{itemize}
    \item A roughly bell-shaped (but positively skewed) distribution
    \item Peak around \$105-\$110
    \item A long right tail extending to \$170+
    \item A sharp left boundary approaching zero (but never reaching it)
    \item The mode is slightly left of the mean, which is left of the tail
\end{itemize}

This is precisely the lognormal distribution!

\newpage

\section{Part VI: Mathematical Justification and Validity}

\subsection{Assumptions Underlying the GBM Model}

\subsubsection{Assumption 1: Constant Parameters}

The GBM model assumes that $\mu$ and $\sigma$ are constants over the forecast horizon.

\paragraph{Justification:}
\begin{itemize}
    \item Over short horizons (1 year), these parameters typically change slowly
    \item Expected returns can be estimated from historical data or market expectations
    \item Volatility can be estimated from historical returns or implied volatility
\end{itemize}

\paragraph{Limitations:}
\begin{itemize}
    \item In reality, volatility is \textbf{stochastic} (varies over time) — this is well-documented in financial markets
    \item Expected returns may shift due to changing economic conditions
    \item For long horizons (5+ years), this assumption becomes less realistic
\end{itemize}

\paragraph{Extensions:}
\begin{itemize}
    \item Stochastic Volatility Models (Heston model)
    \item Jump-Diffusion Models (Merton)
    \item Time-varying parameters
\end{itemize}

\subsubsection{Assumption 2: Lognormality of Returns}

The model assumes that log-returns are normally distributed.

\paragraph{Justification:}
\begin{itemize}
    \item The Central Limit Theorem suggests that returns (as sums of many small shocks) should be approximately normal
    \item For many stocks over moderate time horizons, the normal assumption is reasonable
    \item It makes the model mathematically tractable and closed-form solutions exist
\end{itemize}

\paragraph{Reality Check:}
\begin{itemize}
    \item Real stock returns exhibit \textbf{fat tails}: extreme events occur more frequently than the normal distribution predicts
    \item Real stock returns exhibit \textbf{negative skewness}: crashes are larger and more sudden than rallies
    \item The normal assumption underestimates tail risk
\end{itemize}

\paragraph{Empirical Evidence:}
\begin{table}[H]
\centering
\begin{tabular}{|c|c|c|}
\hline
\textbf{Metric} & \textbf{Normal Distribution} & \textbf{Real Returns} \\
\hline
Probability of $> 3\sigma$ move & 0.27\% & 0.5-1\% \\
Probability of $> 5\sigma$ move & $< 0.00001\%$ & 0.01-0.1\% \\
Skewness & 0 & -0.2 to -0.5 \\
Kurtosis & 3 & 4-10 \\
\hline
\end{tabular}
\caption{Comparison of Normal vs. Real Stock Returns}
\end{table}

\subsubsection{Assumption 3: No Arbitrage}

GBM assumes frictionless markets with no arbitrage opportunities.

\paragraph{Justification:}
\begin{itemize}
    \item Markets are reasonably efficient; arbitrage opportunities are quickly exploited
    \item This is the philosophical foundation of modern option pricing (Black-Scholes)
\end{itemize}

\paragraph{Reality:}
\begin{itemize}
    \item Transaction costs, taxes, and regulatory constraints create small arbitrage-like opportunities
    \item Information asymmetries can lead to mispricings
    \item But for liquid, well-traded assets, the no-arbitrage assumption is reasonable
\end{itemize}

\subsubsection{Assumption 4: Continuous Trading}

The continuous-time formulation assumes prices evolve continuously without gaps.

\paragraph{Reality:}
\begin{itemize}
    \item Markets are closed at night and weekends
    \item Gap moves can occur between close and open
    \item For liquid assets, the continuous assumption is a good approximation over daily or longer horizons
    \item For intra-day analysis, gaps must be modeled explicitly
\end{itemize}

\subsection{When is GBM + Monte Carlo Appropriate?}

\subsubsection{Good Use Cases}

\begin{itemize}
    \item \textbf{Portfolio risk assessment:} Understanding the range of possible portfolio values
    \item \textbf{European-style derivatives:} Option pricing for non-path-dependent options
    \item \textbf{Risk budgeting:} Allocating risk across assets
    \item \textbf{Strategic planning:} Long-term corporate finance decisions
    \item \textbf{Regulatory capital:} Computing Value at Risk for capital adequacy
\end{itemize}

\subsubsection{Poor Use Cases}

\begin{itemize}
    \item \textbf{High-frequency trading:} Too many microstructural effects not captured
    \item \textbf{Extreme tail events:} The normal assumption underestimates rare events
    \item \textbf{Regime shifts:} When $\mu$ and $\sigma$ change dramatically (e.g., during crises)
    \item \textbf{Path-dependent exotics:} Better approached with specific path-dependent models
    \item \textbf{Options on extremes:} Deep out-of-the-money options (tail risk underestimated)
\end{itemize}

\subsection{Model Extensions and Improvements}

\subsubsection{Jump-Diffusion Models}

Extend GBM to include sudden jumps:
\begin{equation}
dS = \mu S \, dt + \sigma S \, dW_t + (J-1) S \, dN_t
\end{equation}

where $dN_t$ is a Poisson process and $J$ is the jump magnitude.

\subsubsection{Stochastic Volatility Models}

Allow volatility to vary over time:
\begin{align}
dS &= \mu S \, dt + \sqrt{v} S \, dW_t^S \\
dv &= \kappa(\theta - v) dt + \xi \sqrt{v} \, dW_t^v
\end{align}

This is the Heston model, which captures volatility clustering.

\subsubsection{Local Volatility Models}

Make volatility a function of current price and time:
\begin{equation}
dS = \mu S \, dt + \sigma(S, t) S \, dW_t
\end{equation}

\subsubsection{Student-t Distributions}

Replace normal distributions with heavier-tailed distributions to capture fat tails.

\newpage

\section{Part VII: Practical Implementation Considerations}

\subsection{Parameter Estimation}

\subsubsection{Estimating $\mu$ (Expected Return)}

\paragraph{Historical Approach:}
\begin{equation}
\hat{\mu} = \frac{1}{n} \sum_{i=1}^{n} r_i
\end{equation}

where $r_i$ are historical log-returns.

\paragraph{Market-Based Approach:}
\begin{itemize}
    \item Use the risk-free rate plus an equity risk premium
    \item Risk-free rate: yield on government bonds
    \item Equity risk premium: historical average 4-8\%, varies by market
\end{itemize}

\paragraph{Analyst Consensus:}
\begin{itemize}
    \item Use forecasts from financial analysts
    \item May incorporate non-historical information
\end{itemize}

\subsubsection{Estimating $\sigma$ (Volatility)}

\paragraph{Historical Volatility:}
\begin{equation}
\hat{\sigma} = \sqrt{\frac{1}{n-1} \sum_{i=1}^{n} (r_i - \bar{r})^2}
\end{equation}

where $r_i$ are log-returns and $\bar{r}$ is the mean.

\paragraph{Implied Volatility:}
\begin{itemize}
    \item Back out volatility from option prices using Black-Scholes
    \item More forward-looking than historical volatility
    \item Reflects market expectations
\end{itemize}

\paragraph{Volatility Decay/Convergence:}
\begin{itemize}
    \item Short-term: Recent volatility matters more
    \item Long-term: Volatility tends to revert to mean
    \item Consider using exponentially weighted moving average (EWMA)
\end{itemize}

\subsection{Variance Reduction Techniques}

To reduce the variance of Monte Carlo estimates without increasing sample size:

\subsubsection{Antithetic Sampling}

For each random draw $Z_i$, also use $-Z_i$. This creates paired simulations that are negatively correlated, reducing variance.

\paragraph{Implementation:}
\begin{equation}
\hat{\mu}_{\text{antithetic}} = \frac{1}{2N} \sum_{i=1}^{N} \left[ f(Z_i) + f(-Z_i) \right]
\end{equation}

\paragraph{Advantage:} Reduces variance by roughly 2x without doubling computation.

\subsubsection{Control Variates}

Use a correlated variable with known expectation to reduce variance.

\paragraph{Idea:} If we know $\mathbb{E}[Y]$ and $Y$ is correlated with our target $X$:
\begin{equation}
\hat{X}_{\text{CV}} = \hat{X} - \beta(\hat{Y} - \mathbb{E}[Y])
\end{equation}

where $\beta = \text{Cov}(X, Y) / \text{Var}(Y)$.

\subsubsection{Importance Sampling}

Sample more frequently from important regions (e.g., tail regions for VaR estimation).

\subsection{Computational Efficiency}

\subsubsection{Parallelization}

Monte Carlo is embarrassingly parallel: each simulation path is independent. 
\begin{itemize}
    \item Each core can independently generate paths
    \item Results are aggregated at the end
    \item Near-linear speedup with number of cores
\end{itemize}

\subsubsection{GPU Acceleration}

GPU (Graphics Processing Unit) computation is particularly well-suited:
\begin{itemize}
    \item Many simple operations (log, exp, sqrt) on different data
    \item GPUs excel at this SIMD (Single Instruction, Multiple Data) pattern
    \item Speedups of 50-500x possible compared to single-core CPU
\end{itemize}

\subsubsection{Adaptive Sampling}

\begin{itemize}
    \item Start with fewer simulations, estimate required precision
    \item Increase sample size adaptively until precision goal is met
    \item Avoid unnecessary over-sampling
\end{itemize}

\subsection{Validation and Testing}

\subsubsection{Backtesting}

Compare historical outcomes to historical simulations:
\begin{itemize}
    \item Simulate over historical periods
    \item Compare empirical distribution to actual outcomes
    \item Check if realized values are within predicted confidence intervals
\end{itemize}

\subsubsection{Stress Testing}

Run simulations under stressed parameter assumptions:
\begin{itemize}
    \item High volatility scenarios
    \item Negative drift scenarios (market downturns)
    \item Correlated asset scenarios
\end{itemize}

\subsubsection{Sensitivity Analysis}

Vary input parameters and observe impact on outputs:
\begin{itemize}
    \item Which parameters most affect the results?
    \item Are parameter estimates robust?
    \item What are the key drivers of risk?
\end{itemize}

\newpage

\section{Part VIII: The Problem → Solution Arc}

\subsection{Recapitulation: The Problem}

\begin{tcolorbox}[colback=red!10, colframe=red!50, title=The Core Problem]

\textbf{Stated Problem:} ``I need to forecast a stock price. But I know it's uncertain. How do I understand that uncertainty quantitatively?''

\textbf{Detailed Problem:}
\begin{enumerate}
    \item \textbf{Drift:} Stocks tend to appreciate, but at an uncertain rate $\mu$
    \item \textbf{Volatility:} Prices fluctuate randomly with standard deviation $\sigma$
    \item \textbf{Proportionality:} Random shocks affect expensive stocks differently than cheap stocks
    \item \textbf{Non-Negativity:} Prices cannot go negative, but a naive model might predict it
    \item \textbf{Asymmetry:} Gains and losses are not symmetric in absolute terms
    \item \textbf{Tail Risk:} I need to understand extreme scenarios, not just the average
    \item \textbf{Correlation:} In a portfolio, I need joint distributions, not just univariate
    \item \textbf{Computation:} I need something numerically tractable for real-world decisions
\end{enumerate}

\end{tcolorbox}

\subsection{How GBM Solves the Problem}

\subsubsection{Problem 1 \& 2: Modeling Drift and Volatility}

\textbf{GBM Solution:}
\begin{equation}
dS = \mu S \, dt + \sigma S \, dW_t
\end{equation}

\begin{itemize}
    \item Drift term $\mu S \, dt$: Captures the expected upward trend
    \item Volatility term $\sigma S \, dW_t$: Captures random fluctuations
    \item Both are built into a single, elegant stochastic differential equation
\end{itemize}

\subsubsection{Problem 3: Proportionality}

\textbf{GBM Solution:} By writing both drift and volatility \textit{multiplicatively} (as fractions of $S$), the model naturally handles proportional shocks.
\begin{itemize}
    \item A $\$10$ drift on a $\$5$ stock is 200\% annual return → unrealistic
    \item A $\$10$ drift on a $\$5000$ stock is 0.2\% annual return → reasonable
    \item The model specifies percentage terms, fixing this implicitly
\end{itemize}

\subsubsection{Problem 4: Non-Negativity}

\textbf{GBM Solution:} The exact solution is:
\begin{equation}
S_t = S_0 \exp\left(\left(\mu - \frac{\sigma^2}{2}\right)t + \sigma W_t\right)
\end{equation}

The exponential function is strictly positive, so $S_t > 0$ with probability 1, regardless of parameter values.

\subsubsection{Problem 5: Asymmetry}

\textbf{GBM Solution:} The exact solution has a lognormal distribution, which is inherently right-skewed. This captures the asymmetry:
\begin{itemize}
    \item Maximum loss: $100\%$ (lose all capital)
    \item Maximum gain: unlimited (theoretically)
    \item The distribution reflects this asymmetry
\end{itemize}

\subsubsection{Problem 6: Tail Risk}

\textbf{GBM + Monte Carlo Solution:}
\begin{itemize}
    \item Run 100,000 simulations
    \item Examine the empirical distribution
    \item Look at the 1st, 5th, and 10th percentiles to understand tail behavior
    \item Compute specific tail metrics (VaR, CVaR)
    \item All extreme scenarios are represented in the simulation distribution
\end{itemize}

\subsubsection{Problem 7: Correlation (Multivariate)}

\textbf{GBM Extension to Multivariate:}

For multiple correlated assets, the multivariate GBM is:
\begin{equation}
dS_i = \mu_i S_i \, dt + \sigma_i S_i \, dW_i^{(1)} + \cdots + \sigma_{i,M} S_i \, dW_i^{(M)}
\end{equation}

where:
\begin{itemize}
    \item Correlations are built in through the common Wiener processes $dW_i^{(j)}$
    \item A correlation matrix $\rho$ can be specified
    \item The simulation generates correlated price paths automatically
\end{itemize}

\subsubsection{Problem 8: Computational Tractability}

\textbf{Monte Carlo Solution:}
\begin{itemize}
    \item The discretized form is simple: $S_{t+\Delta t} = S_t \exp(\ldots)$
    \item Only requires: random normal draws, exponentiations, and multiplications
    \item Easily parallelizable to thousands of cores
    \item Can handle thousands of assets and millions of simulations
    \item Can extend to complex payoff structures and constraints
\end{itemize}

\subsection{Summary: The Complete Solution}

\begin{tcolorbox}[colback=green!10, colframe=green!50, title=The Complete Solution]

\textbf{Model Choice:} Geometric Brownian Motion

\textbf{Why GBM?}
\begin{itemize}
    \item Mathematically elegant and justified by economic theory
    \item Natural non-negativity constraint
    \item Captures both drift and volatility
    \item Produces lognormally-distributed prices (realistic)
    \item Extensible to multivariate and path-dependent scenarios
\end{itemize}

\textbf{Computational Method:} Monte Carlo Simulation

\textbf{Why Monte Carlo?}
\begin{itemize}
    \item Generates empirical distributions reflecting all sources of uncertainty
    \item Simple, interpretable algorithm
    \item Parallelizable for large-scale applications
    \item No dimensional curse: scales reasonably with number of assets
    \item Can incorporate complex features (constraints, exotic payoffs)
\end{itemize}

\textbf{Result:}
\begin{itemize}
    \item A distribution of possible future prices (not a point forecast)
    \item Probabilities of specific outcomes
    \item Risk metrics (VaR, CVaR, volatility)
    \item Insights for decision-making under uncertainty
\end{itemize}

\end{tcolorbox}

\newpage

\section{Conclusion}

Monte Carlo Simulation combined with Geometric Brownian Motion represents a powerful framework for financial forecasting and risk assessment. By grounding ourselves in rigorous stochastic calculus (Itô's Lemma, SDEs), we understand why the model takes the form it does. By discretizing the continuous model, we obtain a computationally tractable algorithm that, when repeated thousands of times, reveals the full landscape of possible futures.

The beauty of this approach is that it:

\begin{enumerate}
    \item \textbf{Respects mathematical constraints} (prices stay positive)
    \item \textbf{Encodes real-world economics} (drift and volatility)
    \item \textbf{Produces realistic distributions} (lognormality, fat tails)
    \item \textbf{Enables quantitative risk assessment} (VaR, probabilities)
    \item \textbf{Is computationally efficient} (parallelizable, scalable)
    \item \textbf{Is theoretically justified} (stochastic calculus, measure theory)
\end{enumerate}

While not perfect—real markets exhibit fat tails, volatility clustering, and regime shifts not captured by the basic GBM—the model provides an excellent first approximation and forms the foundation upon which more sophisticated models are built.

For practitioners, understanding GBM + Monte Carlo is essential. It is the lingua franca of quantitative finance, underpinning option pricing, risk management, portfolio optimization, and strategic financial planning.

\newpage

\appendix

\section{Mathematical Reference}

\subsection{Stochastic Calculus Basics}

\subsubsection{Wiener Process (Brownian Motion)}

A Wiener process $W_t$ satisfies:
\begin{align}
W_0 &= 0 \\
W_t - W_s &\sim \mathcal{N}(0, t-s) \quad \text{for} \quad t > s \\
\text{Cov}(W_s, W_t) &= \min(s, t)
\end{align}

\subsubsection{Itô's Lemma (Stochastic Chain Rule)}

If $dX_t = \mu_X dt + \sigma_X dW_t$ and $Y_t = f(X_t, t)$, then:
\begin{equation}
dY_t = \left(\frac{\partial f}{\partial t} + \mu_X \frac{\partial f}{\partial X} + \frac{1}{2} \sigma_X^2 \frac{\partial^2 f}{\partial X^2}\right) dt + \sigma_X \frac{\partial f}{\partial X} dW_t
\end{equation}

Key difference from deterministic calculus: the $\frac{1}{2} \sigma_X^2 \frac{\partial^2 f}{\partial X^2}$ term is the ``Itô correction.''

\subsubsection{Quadratic Variation}

$(dW_t)^2 = dt$ (not zero, as in deterministic calculus!)

\subsection{Probability Distributions}

\subsubsection{Normal Distribution}

If $X \sim \mathcal{N}(\mu, \sigma^2)$, then:
\begin{align}
\text{PDF}: & \quad f(x) = \frac{1}{\sigma\sqrt{2\pi}} \exp\left(-\frac{(x-\mu)^2}{2\sigma^2}\right) \\
\text{CDF}: & \quad F(x) = \Phi\left(\frac{x-\mu}{\sigma}\right) \\
\text{Mean}: & \quad \mathbb{E}[X] = \mu \\
\text{Variance}: & \quad \text{Var}(X) = \sigma^2
\end{align}

\subsubsection{Lognormal Distribution}

If $X \sim \mathcal{N}(\mu, \sigma^2)$, then $Y = e^X \sim \text{Lognormal}(\mu, \sigma^2)$:
\begin{align}
\text{Mean}: & \quad \mathbb{E}[Y] = \exp\left(\mu + \frac{\sigma^2}{2}\right) \\
\text{Variance}: & \quad \text{Var}(Y) = \left(e^{\sigma^2} - 1\right) e^{2\mu + \sigma^2} \\
\text{Median}: & \quad e^{\mu} \\
\text{Mode}: & \quad e^{\mu - \sigma^2}
\end{align}

\subsection{Central Limit Theorem}

If $X_1, X_2, \ldots$ are i.i.d. with mean $\mu$ and finite variance $\sigma^2$, then:
\begin{equation}
\frac{\sqrt{n}(\bar{X}_n - \mu)}{\sigma} \xrightarrow{d} \mathcal{N}(0, 1)
\end{equation}

This is why sums of many independent random shocks tend to be normally distributed.

\subsection{Law of Large Numbers}

If $X_1, X_2, \ldots$ are i.i.d. with mean $\mu$, then:
\begin{equation}
\bar{X}_n = \frac{1}{n}\sum_{i=1}^n X_i \xrightarrow{p} \mu \quad \text{as } n \to \infty
\end{equation}

This justifies Monte Carlo: the empirical average converges to the true expectation.

\section{Additional Resources and Further Reading}

\begin{itemize}
    \item Hull, J. C. (2021). \textit{Options, Futures, and Other Derivatives} (11th ed.). Pearson.
    \item Björk, T. (2009). \textit{Arbitrage Theory in Continuous Time} (3rd ed.). Oxford University Press.
    \item Glasserman, P. (2004). \textit{Monte Carlo Methods in Financial Engineering}. Springer.
    \item Black, F., \& Scholes, M. (1973). The pricing of options and corporate liabilities. \textit{Journal of Political Economy}, 81(3), 637–654.
    \item Shreve, S. E. (2004). \textit{Stochastic Calculus for Finance} (Vols. I \& II). Springer.
\end{itemize}

\end{document}
